% GENERATED FILE. Edit canonical lessons, then run npm run generate:books.
% canonical-source-hash: fnv1a64:c40562ee0edea90a
% canonical-lessons: GU-C280-jhadpi, GU-W49-jha, GU-C280-ek-tikit, GU-C280-phariyad, GU-C280-khovai-gai, GU-R280-at-the-station

\chapter{Tickets, and when things go wrong}
\label{ch:gu-tickets-and-when-things-go-wrong}
\hlchaptermodality{280}

\begin{chapteropening}
\textbf{By the end of this chapter:} \emph{I can buy a ticket and ask about the next train (āvtī ṭren kyāre che?), report a lost bag or a missed train (khovāī gaī, chūṭī gaī), make a complaint, and read and write \gu{ઝ} in \gu{ઝડપી}.}

Complete GU-R280-at-the-station: buy a ticket, and report a problem.
\end{chapteropening}

\section[\texorpdfstring{\gu{ઝડપી} (jhaḍpī)}{ઝડપી (jhaḍpī)}]{\gu{ઝડપી} (jhaḍpī) — fast; quick}

\label{lesson:GU-C280-jhadpi}



\begin{quote}
Before the new one: ask how to get to the station. (\emph{sṭeśan kevī rīte javũ?}) Then say \textquotedblleft{}a train\textquotedblright{} and \textquotedblleft{}expensive\textquotedblright{}. (\textbf{\gu{ટ્રેન}}, \textbf{\gu{મોંઘું}}.)
\end{quote}

\subsection*{You'll want to know: \gu{ઝડપી}}
\textbf{\gu{ઝડપી}} — \emph{jhaḍpī} — \textquotedblleft{}fast; quick\textquotedblright{}. It ends in \textbf{-\gu{ી}} whatever it describes: \textbf{\gu{ઝડપી} \gu{ટ્રેન}}, \textbf{\gu{ઝડપી} \gu{બસ}}, \textbf{\gu{ઝડપી} \gu{વિમાન}}.

\begin{quote}
\textbf{\gu{ઝડપી} \gu{ટ્રેન} \gu{ક્યારે} \gu{છે}?}
\end{quote}

\begin{quote}
\emph{jhaḍpī ṭren kyāre che?} — \textquotedblleft{}When is the fast train?\textquotedblright{}
\end{quote}

\begin{quote}
\textbf{\gu{તે} \gu{ઝડપી} \gu{છે}, \gu{પણ} \gu{મોંઘી} \gu{છે}.}
\end{quote}

\begin{quote}
\emph{te jhaḍpī che, paṇ moṅghī che.} — \textquotedblleft{}It's fast, but it's expensive.\textquotedblright{}
\end{quote}

Said of the train, \textbf{\gu{મોંઘી}} follows \textbf{\gu{ટ્રેન}}; \textbf{\gu{ઝડપી}} stays as it is.

\begin{cousinweb}
\textbf{\gu{ઝડપી}} (\emph{jhaḍpī}) is built from \textbf{\gu{ઝડપ}}, \textquotedblleft{}speed\textquotedblright{}, and the \textbf{-\gu{ી}} that makes a describing word of it. CDIAL lists Gujarati \emph{jhaṛap} under a reconstructed root \emph{jhaṭ}, \textquotedblleft{}sudden movement\textquotedblright{} (CDIAL 5327). It glosses that form \textquotedblleft{}scuffling\textquotedblright{}, so the sense of speed is probably the later one.
\end{cousinweb}

\subsection*{Your turn}
Say these aloud:

\begin{itemize}
\raggedright
  \item \emph{jhaḍpī}, then \emph{jhaḍpī ṭren}
  \item \textquotedblleft{}When is the fast train?\textquotedblright{} — \emph{jhaḍpī ṭren kyāre che?}
  \item \textquotedblleft{}It's fast, but it's expensive\textquotedblright{} — \emph{te jhaḍpī che, paṇ moṅghī che}
\end{itemize}

\begin{tcolorbox}[breakable,colback=teal!4,colframe=teal!35!black,title={Before you move on}]
What does \textbf{\gu{ઝડપી}} mean? (\textquotedblleft{}Fast; quick\textquotedblright{}.) \textquotedblleft{}A fast bus\textquotedblright{}? (\textbf{\gu{ઝડપી} \gu{બસ}}.) Does \textbf{\gu{ઝડપી}} change for a bus? (No.)
\end{tcolorbox}
\section[\texorpdfstring{\gu{ઝ} (jha)}{ઝ (jha)}]{\gu{ઝ} — the consonant jha}

\label{lesson:GU-W49-jha}



\begin{quote}
\begin{itemize}
\raggedright
  \item \emph{Recall:} say \emph{fast} — \textbf{R1}, one lesson back
  \item \emph{Recall:} say \textbf{\gu{જ}}, as in \textbf{\gu{જવું}}, and \textbf{\gu{છ}}, as in \textbf{\gu{છે}}
\end{itemize}

\textbf{\gu{ઝડપી}} opens on a sound you have not written yet: \emph{jh}.
\end{quote}

\subsection*{You'll want to know: \gu{જ} with a breath}
\noindent
\begin{tabularx}{\linewidth}{@{}>{\raggedright\arraybackslash}X>{\raggedright\arraybackslash}X>{\raggedright\arraybackslash}X@{}}
\toprule
\textbf{} & \textbf{plain} & \textbf{with a breath} \\
\midrule
\emph{c} & \textbf{\gu{ચ}} & \textbf{\gu{છ}} \\
\emph{j} & \textbf{\gu{જ}} & ? \\
\bottomrule
\end{tabularx}

The missing letter is \textbf{\gu{ઝ}}, \emph{jha}: say \textbf{\gu{જ}} and let a breath follow. Its print has three parts: a rounded left body, a separate right loop flowing into a lower tail, and a separate short upper stem.

\subsection*{Script — meet one form}
\begin{quote}
\textbf{\gu{ઝ}} — \emph{jha}
\end{quote}

Keep \textbf{\gu{ઝ}} visible. Follow the numbered strip: it shows where each part starts.

\begin{enumerate}
\raggedright
  \item begin at the upper left and circle the rounded left body to the lower left
  \item lift once and circle the right loop into its lower tail
  \item lift again and descend the short upper stem
\end{enumerate}

Three strokes, two pen lifts: a verified three-stroke teaching-app variant fitted to the bundled printed outline (t30apps.com, ‘Gujarati Alphabet Writing Practice’, version 1.0, \gu{ઝ} animation, first through third SVG paths).

\subsection*{Writing — notice and trace}
\hlblockfigure{\detokenize{figures/GU-W49-jha-filmstrip.pdf}}{How \gu{ઝ} is written, stroke by stroke}

Finger-trace \textbf{\gu{ઝ}} once, then air-trace it once. Notice the two lifts: after the left body, and after the loop.

\subsection*{Writing — copy with the model}
Copy visible \textbf{\gu{ઝ}} once. Compare only two things: are there three parts, and does the short stem come last? Repair one landmark if needed.

\subsection*{Writing — recall after hiding}
Hide the model and write \emph{jha} once. Uncover \textbf{\gu{ઝ}}, compare, and stop after one repair.

\begin{tcolorbox}[breakable,colback=teal!4,colframe=teal!35!black,title={Before you move on}]
Read \textbf{\gu{જ} \gu{ઝ} \gu{ઝડપી}} aloud. Which part of \textbf{\gu{ઝ}} is written last? (\textbf{The short upper stem.})

Sources: \href{https://www.t30apps.com/gujarati-alphabet-writing/}{Gujarati Alphabet Writing Practice} and the \href{https://www.unicode.org/charts/PDF/U0A80.pdf}{Unicode Gujarati chart}.
\end{tcolorbox}
\section[\texorpdfstring{\gu{એક} \gu{ટિકિટ} \gu{આપો} (ek ṭikiṭ āpo)}{એક ટિકિટ આપો (ek ṭikiṭ āpo)}]{\gu{એક} \gu{ટિકિટ} \gu{આપો} (ek ṭikiṭ āpo) — at the ticket window}

\label{lesson:GU-C280-ek-tikit}



\begin{quote}
Say \textquotedblleft{}the fast train\textquotedblright{}. (\emph{jhaḍpī ṭren}.) Then \textquotedblleft{}a ticket\textquotedblright{} and \textquotedblleft{}to give\textquotedblright{}. (\textbf{\gu{ટિકિટ}}, \textbf{\gu{આપવું}}.)
\end{quote}

\begin{grammarlens}[title={Grammar Lens: asking for a ticket}]
\begin{quote}
\textbf{\gu{બે} \gu{ટિકિટ} \gu{આપો}.}
\end{quote}

\begin{quote}
\emph{be ṭikiṭ āpo.} — \textquotedblleft{}Two tickets, please.\textquotedblright{}
\end{quote}

\textbf{\gu{આપો}} is \textbf{\gu{આપવું}} in the \textbf{-\gu{ઓ}} form of \textbf{\gu{જાઓ}} and \textbf{\gu{વળો}}: said with respect at a window, it is polite enough. Then the price: \textbf{\gu{ભાડું} \gu{કેટલું} \gu{છે}?} (\emph{bhāḍũ keṭlũ che?}), \textquotedblleft{}how much is the fare?\textquotedblright{} \textbf{\gu{ઝડપી} \gu{ટ્રેનની} \gu{ટિકિટ} \gu{મોંઘી} \gu{છે}} (\emph{jhaḍpī ṭrennī ṭikiṭ moṅghī che}): the fast train's ticket costs more.
\end{grammarlens}

\subsection*{You'll want to know: the next train}
\begin{quote}
\textbf{\gu{આવતી} \gu{ટ્રેન} \gu{ક્યારે} \gu{છે}?}
\end{quote}

\begin{quote}
\emph{āvtī ṭren kyāre che?} — \textquotedblleft{}When is the next train?\textquotedblright{}
\end{quote}

\textbf{\gu{આવતું}}, \textquotedblleft{}coming, next\textquotedblright{}, follows \textbf{\gu{ટ્રેન}}: \textbf{\gu{આવતી}}. Then the platform: \textbf{\gu{ટ્રેન} \gu{કયા} \gu{પ્લેટફોર્મથી} \gu{ઊપડશે}?} (\emph{ṭren kayā pleṭphormthī ūpaḍśe?}), \textquotedblleft{}which platform does the train leave from?\textquotedblright{} The answer: \textbf{\gu{પ્લેટફોર્મ} \gu{બેથી}} (\emph{pleṭphorm bethī}), \textquotedblleft{}from platform two\textquotedblright{}.

\subsection*{Your turn}
Say these aloud:

\begin{itemize}
\raggedright
  \item \textquotedblleft{}How much is the fare?\textquotedblright{} — \emph{bhāḍũ keṭlũ che?}
  \item \textquotedblleft{}When is the next bus?\textquotedblright{} — \emph{āvtī bas kyāre che?}
\end{itemize}

\begin{tcolorbox}[breakable,colback=teal!4,colframe=teal!35!black,title={Before you move on}]
\textquotedblleft{}One ticket, please\textquotedblright{}? (\textbf{\gu{એક} \gu{ટિકિટ} \gu{આપો}}.) \textquotedblleft{}The next train\textquotedblright{}? (\textbf{\gu{આવતી} \gu{ટ્રેન}}.) \textquotedblleft{}From which platform?\textquotedblright{} (\textbf{\gu{કયા} \gu{પ્લેટફોર્મથી}?})
\end{tcolorbox}
\section[\texorpdfstring{\gu{ફરિયાદ} (phariyād)}{ફરિયાદ (phariyād)}]{\gu{ફરિયાદ} (phariyād) — a complaint}

\label{lesson:GU-C280-phariyad}



\begin{quote}
Before the new one: ask when the next train is. (\emph{āvtī ṭren kyāre che?}) Then say \textquotedblleft{}the police\textquotedblright{}. (\textbf{\gu{પોલીસ}}.)
\end{quote}

\subsection*{You'll want to know: \gu{ફરિયાદ}}
\textbf{\gu{ફરિયાદ}} — \emph{phariyād} — \textquotedblleft{}a complaint\textquotedblright{}. It is feminine, so a past of \textbf{\gu{કરવું}} follows it with \textbf{\gu{કરી}}:

\begin{quote}
\textbf{\gu{મેં} \gu{પોલીસમાં} \gu{ફરિયાદ} \gu{કરી}.}
\end{quote}

\begin{quote}
\emph{mẽ polīsmā̃ phariyād karī.} — \textquotedblleft{}I made a complaint to the police.\textquotedblright{}
\end{quote}

\textbf{\gu{પોલીસમાં}} (\emph{polīsmā̃}) is \textquotedblleft{}with the police\textquotedblright{}, word for word \textquotedblleft{}in the police\textquotedblright{}. A complaint can be about anything:

\begin{quote}
\textbf{\gu{તમારી} \gu{ફરિયાદ} \gu{શું} \gu{છે}?}
\end{quote}

\begin{quote}
\emph{tamārī phariyād śũ che?} — \textquotedblleft{}What is your complaint?\textquotedblright{}
\end{quote}

\begin{cousinweb}
\textbf{\gu{ફરિયાદ}} (\emph{phariyād}) is a loan, usually traced to Persian \emph{faryād}, \textquotedblleft{}an outcry; a call for help\textquotedblright{}, and from there \textquotedblleft{}a complaint\textquotedblright{}.
\end{cousinweb}

\subsection*{Your turn}
Say these aloud:

\begin{itemize}
\raggedright
  \item \emph{phariyād}
  \item \textquotedblleft{}I made a complaint to the police\textquotedblright{} — \emph{mẽ polīsmā̃ phariyād karī}
  \item \textquotedblleft{}What is your complaint?\textquotedblright{} — \emph{tamārī phariyād śũ che?}
\end{itemize}

\begin{tcolorbox}[breakable,colback=teal!4,colframe=teal!35!black,title={Before you move on}]
What does \textbf{\gu{ફરિયાદ}} mean? (\textquotedblleft{}A complaint\textquotedblright{}.) Why \textbf{\gu{કરી}} in \textbf{\gu{મેં} \gu{ફરિયાદ} \gu{કરી}}? (\textbf{\gu{ફરિયાદ}} is feminine.)
\end{tcolorbox}
\section[\texorpdfstring{\gu{મારી} \gu{થેલી} \gu{ખોવાઈ} \gu{ગઈ} (mārī thelī …)}{મારી થેલી ખોવાઈ ગઈ (mārī thelī …)}]{\gu{મારી} \gu{થેલી} \gu{ખોવાઈ} \gu{ગઈ} (mārī thelī khovāī gaī) — when things go wrong}

\label{lesson:GU-C280-khovai-gai}



\begin{quote}
Say \textquotedblleft{}a complaint\textquotedblright{}. (\emph{phariyād}.) Then \textquotedblleft{}to get lost\textquotedblright{} and \textquotedblleft{}a bag\textquotedblright{}. (\textbf{\gu{ખોવાવું}}, \textbf{\gu{થેલી}}.)
\end{quote}

\begin{grammarlens}[title={Grammar Lens: the lost thing is the subject}]
\begin{quote}
\textbf{\gu{મારી} \gu{થેલી} \gu{ખોવાઈ} \gu{ગઈ}.}
\end{quote}

\begin{quote}
\emph{mārī thelī khovāī gaī.} — \textquotedblleft{}My bag is lost.\textquotedblright{}
\end{quote}

The bag is what got lost, so \textbf{\gu{ગઈ}} follows it, and there is no \textbf{-\gu{એ}}. Change the thing and the ending changes: \textbf{\gu{મારું} \gu{પાકીટ} \gu{ખોવાઈ} \gu{ગયું}} (\emph{mārũ pākīṭ khovāī gayũ}), neuter; \textbf{\gu{મારો} \gu{ફોન} \gu{ખોવાઈ} \gu{ગયો}} (\emph{māro phon khovāī gayo}), masculine.
\end{grammarlens}

\subsection*{You'll want to know: a missed train, and help}
\begin{quote}
\textbf{\gu{મારી} \gu{ટ્રેન} \gu{છૂટી} \gu{ગઈ}.}
\end{quote}

\begin{quote}
\emph{mārī ṭren chūṭī gaī.} — \textquotedblleft{}I missed my train.\textquotedblright{}
\end{quote}

Word for word, \textquotedblleft{}my train got away\textquotedblright{}: the train is the subject again. Then ask for help: \textbf{\gu{મને} \gu{મદદ} \gu{કરશો}?} (\emph{mane madad karśo?}), \textquotedblleft{}will you help me?\textquotedblright{} And the next step: \textbf{\gu{મારે} \gu{પોલીસમાં} \gu{ફરિયાદ} \gu{કરવી} \gu{પડશે}} (\emph{māre polīsmā̃ phariyād karvī paḍśe}), \textquotedblleft{}I'll have to make a complaint to the police\textquotedblright{}. \textbf{\gu{કરવી}} follows \textbf{\gu{ફરિયાદ}}.

\subsection*{Your turn}
Say these aloud:

\begin{itemize}
\raggedright
  \item \textquotedblleft{}My key is lost\textquotedblright{} — \emph{mārī cāvī khovāī gaī}
  \item \textquotedblleft{}I missed my train, so I'll have to take the bus\textquotedblright{} — \emph{mārī ṭren chūṭī gaī, eṭle māre basmā̃ javũ paḍśe}
\end{itemize}

\begin{tcolorbox}[breakable,colback=teal!4,colframe=teal!35!black,title={Before you move on}]
\textquotedblleft{}My bag is lost\textquotedblright{}? (\textbf{\gu{મારી} \gu{થેલી} \gu{ખોવાઈ} \gu{ગઈ}}.) \textquotedblleft{}My phone is lost\textquotedblright{}? (\textbf{\gu{મારો} \gu{ફોન} \gu{ખોવાઈ} \gu{ગયો}}.) \textquotedblleft{}I missed my train\textquotedblright{}? (\textbf{\gu{મારી} \gu{ટ્રેન} \gu{છૂટી} \gu{ગઈ}}.)
\end{tcolorbox}
\section[(review)]{Review — at the station}

\label{lesson:GU-R280-at-the-station}



\begin{quote}
Ask \textquotedblleft{}how do I get to the station?\textquotedblright{} (\emph{sṭeśan kevī rīte javũ?}) Then \textquotedblleft{}turn left\textquotedblright{}. (\emph{ḍābī bāju vaḷo}.)

\begin{itemize}
\raggedright
  \item \emph{Write it:} \textbf{\gu{ઝ}}, then \textbf{\gu{ઝડપી}}, from memory
\end{itemize}
\end{quote}

\subsection*{Your turn}
Say these aloud:

\begin{itemize}
\raggedright
  \item Two tickets, please. How much is the fare?
  \item When is the next train? Which platform does it leave from?
  \item The fast train is quick, but it's expensive.
  \item My bag is lost. My wallet is lost too.
  \item I missed my train. Will you help me? I'll have to make a complaint to the police.
\end{itemize}

\begin{tcolorbox}[breakable,colback=teal!4,colframe=teal!35!black,title={Before you move on}]
\textquotedblleft{}The next train\textquotedblright{}? (\textbf{\gu{આવતી} \gu{ટ્રેન}}.) \textquotedblleft{}My phone is lost\textquotedblright{}: \textbf{\gu{ગયો}} or \textbf{\gu{ગઈ}}? (\textbf{\gu{ગયો}}: \textbf{\gu{ફોન}} is masculine.) Which letter opens \textbf{\gu{ઝડપી}}? (\textbf{\gu{ઝ}}, \emph{jha}.)
\end{tcolorbox}
